%% ma: L12-B08-0108
%% lop: 12
%% chuong: 2
%% bai: 8
%% dang: 12.8.7
%% loai: TN4
%% muc: NB
%% dap_an: "B"
%% nguon: "CĐ2.3.pdf giáo viên gửi 27/10/2026 (M07, Lê Minh Khôi – Nguyễn Anh Khôi, Chương II Chủ đề 3 Biểu thức tọa độ của các phép toán vectơ), Câu 19"
%% kien_thuc: "Bài 8 (tọa độ trung điểm)"
%% trang_thai: chinh-thuc
\begin{ex}
Trong không gian $Oxyz$, cho hai điểm $A(1;-3;2)$, $B(3;-1;4)$. Tọa độ trung điểm $I$ của đoạn thẳng $AB$ là
\choice{$I(2;2;2)$}{\True $I(2;-2;3)$}{$I(1;1;1)$}{$I(4;-4;6)$}
\loigiai{$I\left(\dfrac{1+3}2;\dfrac{-3-1}2;\dfrac{2+4}2\right)=(2;-2;3)$.}
\end{ex}
